\chapter{Vaginal thrush}
\index{Vaginal thrush}

\section*{Definition and Overview}
Vaginal thrush, also known as vulvovaginal candidiasis, is a very common fungal infection of the vagina and the surrounding vulval skin. It is caused by an overgrowth of yeast of the genus \textit{Candida}, most often \textit{Candida albicans}. Thrush produces a classic set of symptoms: intense vulval itching and soreness, a thick white discharge that is often described as resembling cottage cheese, redness, and a burning sensation, particularly with urination or intercourse.

Thrush is not a sexually transmitted infection, and it can affect women who have never been sexually active. It is usually easy to treat with antifungal creams, pessaries, or tablets, but it tends to recur, and a small proportion of women experience frequent or resistant episodes that require specialist help. Because other conditions produce similar symptoms, it is wise to have an accurate diagnosis, especially when symptoms are recurrent, unusual, or not settling with over-the-counter treatment.

\section*{Epidemiology}
Vaginal thrush is one of the most common infections seen in primary care worldwide. Around 75 percent of women experience at least one episode of vulvovaginal candidiasis during their lifetime, and roughly 40 to 50 percent have two or more attacks. Uncomplicated thrush is most common in women of childbearing age, and its frequency peaks when hormonal influences are strongest, particularly during the second trimester of pregnancy.

Approximately 5 to 8 percent of women suffer from recurrent vulvovaginal candidiasis, usually defined as four or more confirmed episodes in a single year. Certain groups are at higher risk: women taking broad-spectrum antibiotics, women with diabetes, particularly poorly controlled disease, pregnant women, and women taking combined oral contraceptives or undergoing hormone replacement. Thrush is also more common and often more severe in women with weakened immunity, including those living with HIV, though in otherwise healthy women severe thrush is uncommon and should prompt consideration of an underlying cause such as diabetes.

\section*{Aetiology and Pathophysiology}
\textit{Candida} yeast is a normal, harmless inhabitant of the skin, bowel, and vagina in many healthy women. Disease occurs when the balance of the vaginal ecosystem tips in favour of yeast overgrowth. The most important protective mechanisms are the acidic pH maintained by \textit{Lactobacillus} bacteria and the physical barrier of the vaginal lining; anything that disturbs these can allow \textit{Candida} to multiply and change from a budding yeast into the invasive, thread-like hyphal form that damages tissue and triggers inflammation.

\subsection*{Predisposing factors}
\begin{itemize}
    \item \textbf{Antibiotics:} kill the protective \textit{Lactobacillus} bacteria, allowing yeast to flourish
    \item \textbf{Hormonal change:} high oestrogen states such as pregnancy and the use of combined hormonal contraceptives favour yeast growth
    \item \textbf{Diabetes:} high sugar levels in vaginal secretions feed the yeast
    \item \textbf{Immunosuppression:} HIV, chemotherapy, long-term corticosteroids, and other causes of reduced immunity
    \item \textbf{Lifestyle factors:} tight synthetic clothing, damp underwear, and heat may create a moist environment that encourages growth
    \item \textbf{Skin damage:} eczema, dermatitis, or irritation from soaps and products can disrupt the protective surface
\end{itemize}

\section*{Clinical Features}
\begin{itemize}
    \item \textbf{Itching:} intense vulval itching is the hallmark symptom and may be worse at night
    \item \textbf{Discharge:} thick, white, lumpy discharge resembling cottage cheese, although some women have only a scant or watery discharge
    \item \textbf{Soreness and redness:} the vulva and vagina are red, swollen, and tender; tiny cracks (fissures) may develop in the skin
    \item \textbf{Burning:} burning with urination or during and after intercourse
    \item \textbf{Rash:} small satellite spots of dermatitis may extend beyond the main red area onto the thighs
    \item \textbf{Absent signs:} some women, particularly with recurrent disease, have surprisingly few visible changes despite severe symptoms
\end{itemize}

The intensity of symptoms varies widely. At one end are women with mild, occasional itching and a barely noticeable change in discharge; at the other are women with excoriated, raw vulval skin who find sitting, walking, and urinating painful. Symptoms typically flare in the week before a period, and the first episode in a woman's life is often the most severe because the skin has not previously been sensitised. Because itching is the dominant symptom, women often scratch, which worsens the redness, swelling, and fissuring and can allow bacteria to enter. Keeping the area cool, dry, and free of soap and fragrance is helpful from the moment symptoms begin, even before treatment starts.

\section*{Differential Diagnosis}
\begin{itemize}
    \item Bacterial vaginosis (thin, fishy-smelling discharge with little itching)
    \item Trichomoniasis (frothy, yellow-green, offensive discharge from a sexually transmitted parasite)
    \item Chlamydia and gonorrhoea (often asymptomatic but can cause discharge and irritation)
    \item Contact dermatitis or allergic reactions to soaps, detergents, and intimate products
    \item Vulval skin conditions such as lichen sclerosus and eczema
    \item Atrophic vaginitis in older women, where dryness and soreness predominate
    \item Vulvodynia, a chronic pain disorder without visible abnormality
\end{itemize}

\section*{When to Seek Medical Care}
See a doctor if this is your first episode of thrush-like symptoms, if you are pregnant, if symptoms do not settle after completing an over-the-counter course of treatment, if discharge becomes blood-stained, or if you have diabetes, a weakened immune system, or frequent recurrences. A swab test can confirm the diagnosis and, importantly, exclude other infections that need different treatment.

\textbf{Contact your local emergency services for:}
\begin{itemize}
    \item Severe lower abdominal or pelvic pain, especially with fever
    \item Heavy vaginal bleeding or signs of a severe general illness
    \item Redness spreading rapidly with severe pain and swelling, which can indicate a bacterial skin infection
    \item Any sudden, severe, or concerning symptom
\end{itemize}

\section*{Diagnosis and Investigations}
\begin{itemize}
    \item \textbf{History:} characteristic symptoms, recent antibiotic use, diabetes, pregnancy, and previous episodes
    \item \textbf{Examination:} inspection of the vulva and vagina shows redness, swelling, fissures, and typical discharge
    \item \textbf{Microscopy:} a swab examined under the microscope reveals budding yeast cells and hyphae
    \item \textbf{PH testing:} vaginal pH is usually normal (4.5 or below) in thrush, which helps distinguish it from bacterial vaginosis
    \item \textbf{Culture:} growing the yeast in the laboratory identifies the species, which matters if the infection is resistant; culture is often used in recurrent disease
    \item \textbf{Blood tests:} blood sugar testing is recommended if there is a suggestion of diabetes
    \item \textbf{Specialist review:} a gynaecologist or infectious-disease specialist may manage severe, recurrent, or resistant cases
\end{itemize}

\section*{Management}
\subsection*{First-line treatment of uncomplicated thrush}
\begin{itemize}
    \item \textbf{Antifungal creams:} clotrimazole or miconazole applied to the vulva, usually for one to two weeks
    \item \textbf{Vaginal pessaries:} clotrimazole or miconazole inserted into the vagina, sometimes as a single high dose
    \item \textbf{Oral tablet:} fluconazole as a single dose is effective, convenient, and popular; a doctor may prescribe it for those who prefer not to use vaginal products
\end{itemize}

\subsection*{Recurrent or complicated thrush}
Recurrent thrush (four or more episodes a year) is treated with induction treatment to clear the current attack followed by a maintenance regimen, often a regular oral or vaginal antifungal taken weekly for several months. Underlying factors such as diabetes are addressed, and tests may be arranged to confirm the diagnosis at each episode, since many presumed "recurrences" are actually bacterial vaginosis or other conditions. \textit{Candida glabrata} and other non-\textit{albicans} species are more resistant to standard antifungals and may need alternative agents under specialist supervision.

\subsection*{Pregnancy}
Thrush is very common in pregnancy. Topical antifungals are safe and effective, but oral fluconazole is generally avoided in early pregnancy, so a doctor or midwife should advise on the safest option. Symptoms in pregnancy can be severe, and episodes may need repeated treatment.

\subsection*{General measures}
Avoid soap, perfumed products, and tight clothing while treating an episode; use cool compresses for relief; and complete the full course of treatment. Sexual partners do not require treatment because thrush is not sexually transmitted, although intercourse may be uncomfortable while symptoms persist.

\section*{Complications}
\begin{itemize}
    \item Recurrence, which is the most common problem and can be distressing and disruptive
    \item Severe vulval inflammation, fissuring, and secondary bacterial infection in some cases
    \item Skin thickening and lichenification from persistent scratching
    \item In pregnancy, thrush is uncomfortable but does not harm the baby; rare cases of oral thrush (thrush in the baby's mouth) can occur after birth
    \item Invasive candidiasis is essentially confined to people with severely weakened immune systems and is not a concern for otherwise healthy women
\end{itemize}

\section*{Prognosis and Outcomes}
For most women, an episode of vaginal thrush clears completely within a few days of starting treatment, and oral and topical treatments are similarly effective. Around 40 to 50 percent of women experience a further episode within a year, but for the majority these remain occasional and easily managed. For the smaller group with recurrent disease, maintenance treatment substantially reduces the frequency of attacks while it is continued, although symptoms may return once it is stopped. Thrush does not damage the reproductive organs and has no lasting effect on fertility, pregnancy, or long-term health. Working with a doctor to confirm the diagnosis and address any underlying cause offers the best chance of long-term control.

It is worth emphasising that most women who have a single episode will not go on to have a problem at all. Recurrent thrush is more likely when a trigger such as repeated antibiotics, poorly controlled diabetes, or a persistent moisture problem remains, so identifying and managing the trigger is the most effective long-term strategy. When recurrence does occur, it is important to confirm that it really is thrush, since roughly one in three "recurrent thrush" cases turns out to be another condition such as bacterial vaginosis. With accurate diagnosis, management of triggers, and, where needed, maintenance treatment, even the most troublesome patterns of thrush can usually be brought under control.

\section*{Prevention and Self-Care}
\begin{itemize}
    \item Avoid douching and perfumed soaps, deodorants, and bubble bath
    \item Wear cotton underwear, avoid tight synthetic clothing, and change out of wet swimwear or exercise clothes promptly
    \item Wipe from front to back after using the toilet
    \item If thrush follows antibiotic courses, discuss preventive antifungals with your doctor
    \item Keep diabetes well controlled
    \item Do not use treatments for longer than recommended without medical advice, as unnecessary antifungals can promote resistance
    \item See a doctor to confirm the diagnosis if symptoms recur, rather than repeatedly self-treating
\end{itemize}

It is worth repeating that self-diagnosis is unreliable. Up to half of women who treat presumed thrush with over-the-counter antifungals turn out to have bacterial vaginosis, a mixed infection, or no infection at all, which is why recurrent or atypical symptoms deserve a swab test. A vaginal swab is quick, painless, and taken in the privacy of a consultation or sometimes self-collected. Confirming the organism, checking the vaginal pH, and excluding other infections gives treatment the best chance of working and stops the cycle of repeated courses of the wrong medicine. For most women, however, an occasional typical episode of thrush responds well to a standard antifungal, and the information in this chapter will help them use it correctly and know when to seek help.

\section*{Sources and Further Reading}
The following organisations provide reliable information on vaginal thrush and its treatment.

\begin{itemize}
    \item NHS. \textit{Information on thrush in women, causes, and treatment.} https://www.nhs.uk/conditions/
    \item Mayo Clinic. \textit{Guide to vaginal yeast infection, symptoms, and treatment.} https://www.mayoclinic.org/
    \item MedlinePlus. \textit{Health information on candidiasis and vaginal fungal infections.} https://medlineplus.gov/
    \item MSD Manual. \textit{Patient education on candidiasis and vulvovaginal infections.} https://www.msdmanuals.com/
    \item HealthDirect Australia. \textit{Consumer information on thrush and vaginal infections.} https://www.healthdirect.gov.au/
    \item Jean Hailes for Women's Health. \textit{Resources on vulval and vaginal health.} https://www.jeanhailes.org.au/
\end{itemize}
